\section{Massabielle: Poverty, Publicity, and a Growing Record}

Lourdes in 1858 was a small Pyrenean town within the Diocese of Tarbes and the political order of France's Second Empire. Massabielle lay near the Gave de Pau, outside ordinary respectability. The name's association with an old rock or mass and the grotto's use as rough marginal ground matter symbolically, but the first historical point is simpler: Bernadette went there with her sister and a companion to gather fuel and bones because the household was poor.

The official site now describes the grotto as a limestone cavity about 3.8 metres high, 9.5 metres deep, and 9.85 metres wide, with a karstic extension of roughly ten metres. These measurements belong to present site description, not to a supernatural proof. They establish that the scene was material and local: cold water, rock, mud, vegetation, a river, changing access, and bodies affected by weather and illness.

\subsection{The Soubirous household}

François Soubirous and Louise Castérot had known the relative security of the Boly mill. Economic failure, debt, unstable work, and suspicion drove the family downward. By 1857 they lived in the \emph{Cachot}, a former jail room. Bernadette, born 7 January 1844 and baptized two days later, had survived cholera and lived with chronic respiratory weakness. She spent time at Bartrès helping with animals and returned to Lourdes in January 1858, wanting the catechesis needed for First Communion.

Poverty can expose a witness to contempt, pressure, dependency, and romantic appropriation. It does not make a person naturally more truthful or supernaturally chosen. Nor does limited schooling mean an empty mind. Bernadette knew household Catholicism, elementary prayers, the sign of the Cross, the Rosary, parish relationships, and the imaginative world of rural devotion. Any historical argument from what she “could not possibly know” must therefore be precise rather than theatrical.

The Gospel repeatedly discloses God's freedom to choose the poor. That theological pattern illuminates Lourdes after the evidence is named; it does not replace evidence. The correct reversal is moral: a society that treated the family as disposable was compelled to hear a girl address clergy and officials without becoming their possession.

\subsection{A claim public almost at once}

The Lourdes story did not first emerge decades later. Bernadette told her companions; family and neighbors responded; crowds grew; clergy maintained caution; police and magistrates questioned her; officials attempted to control access; physicians watched her bodily states; and reports multiplied during the sequence itself. On 21 February Commissioner Dominique Jacomet questioned her, and the received chronology preserves her early word \emph{Aquerò}---“that one”---rather than a developed theological identification.

This density is evidentially valuable. Repeated questioning tests memory, consistency, suggestibility, and motive. Publicity also contaminates memory, supplies vocabulary, intensifies expectation, and creates social roles. Near-event abundance is not the same as a controlled experiment. The historian asks who recorded a statement, when, in what language, after what prompting, and how the surviving text was edited.

Crowds likewise prove less and more than devotional shorthand suggests. They prove that Bernadette's behavior was public and that the grotto rapidly became a contested space. They can corroborate her movements, expression, prayer, and the presence of civil authorities. The received corpus does not claim that eight thousand persons on 4 March all saw the Lady. The visionary object remained Bernadette's experience.

\sourceboundary{Bernadette's repeated 1858 statements and observed behavior}{A stable claimant, a bounded sequence, public scrutiny, and a message developing before the final title}{A shared public visual object, dictated transcript, or automatic exclusion of every natural psychological process}

\subsection{The first autograph and its limit}

The Sanctuary publishes an extract from Bernadette's first autograph account, addressed to Father Gondrand of the Oblates of Mary Immaculate at Bétharram. It explicitly notes that the writing followed numerous oral narratives and official “statements.” The manuscript is therefore primary first-person evidence, not a contemporaneous diary entry written at noon on 11 February.

Its concrete restraint is important. Bernadette describes a white gown, white veil, blue sash, golden roses, and a rosary. She first cannot raise her hand for the sign of the Cross; after the Lady signs herself, Bernadette can. Bernadette then recites the Rosary. The figure passes the beads through her fingers \emph{without moving her lips}. This is not a report that the Lady audibly spoke the Hail Mary. Later phrases such as “praying the Rosary with Mary” are legitimate devotional synthesis only if they preserve that sensory claim.

The account also says the two companions noticed nothing. Visionary asymmetry is not an embarrassment to be removed. It identifies the kind of claim under examination: an experience reported by Bernadette, accompanied by outward behavior visible to others, rather than an ordinary object jointly perceived by everyone present.

\subsection{From inquiry to judgment}

On 28 July 1858 Bishop Laurence constituted a commission to collect and establish events at and around the grotto, characterize them, and supply the elements needed for judgment. The roughly three-and-a-half-year interval was not indifference. It allowed the eighteen-event cycle to close, testimony and objections to accumulate, early cures to be examined, and fruits to become more legible.

The bishop's reasoning, as summarized by the Sanctuary, attended to the witness's reliability, spiritual fruits, and bodily cures. Modern readers should neither mock nor idealize that framework. Witness, doctrine, fruit, and physical phenomena remain relevant. Modern inquiry would add developed psychological, safeguarding, evidentiary, and medical protocols. The 1862 act should be assessed as a serious historical ecclesial judgment, not impersonated as a twenty-first-century trial.

The complete Tarbes archive and the full critical documentary edition were not collated for this revision. The official Sanctuary's exact operative transcription, English status page, detailed chronology, compressed calendar, first-autograph extract, and current institutional records form the verified working corpus. This limitation is material. It prevents unlocated deposition details from appearing as though directly checked.

\subsection{Received narrative and honest harmonization}

A shrine must teach a coherent story to pilgrims. It condenses repeated testimonies, later recollections, civil records, art, liturgy, and pastoral interpretation. That is not inherently falsification. Trouble begins when the resulting narrative loses its source labels.

The Lourdes sequence has a remarkably stable spine: two initially wordless encounters; an invitation to return; prayer and penance; sinners; the spring; chapel and procession; a silent conclusion to the fortnight; the title on 25 March; the candle episode; and the final silent vision in July. The placement of a private prayer and secrets varies between two current official summaries. Since those elements were not a public teaching, doctrinal synthesis does not depend upon solving the variation.

Good source criticism thus serves rather than dissolves devotion. It removes imagined blossoms, invented audible prayers, published “secrets,” and automatic cure claims so the actual center can be heard.
